\chapter{راهبرد تعامل بین‌المللی و سیاست خارجی}
\label{app:intl}

ایران نوین نیازمند بازگشت به جامعه‌ی جهانی به عنوان یک بازیگر مسئول، صلح‌طلب و پیش‌بینی‌پذیر است.

\section{اصول راهنمای سیاست خارجی}

\begin{modernbox}{دکترین «ایران برای صلح»}
\begin{enumerate}
    \item \textbf{منافع ملی:} اولویت مطلق بر اساس رفاه و امنیت ملت ایران.
    \item \textbf{تنش‌زدایی منطقه‌ای:} حل اختلافات با همسایگان و پایان دادن به جنگ‌های نیابتی.
    \item \textbf{ادغام اقتصادی:} الحاق به WTO و بازگشت به بازارهای جهانی انرژی.
    \item \textbf{پایبندی به حقوق بشر:} الحاق بدون قید و شرط به کنوانسیون‌های بین‌المللی.
    \item \textbf{شفافیت هسته‌ای:} حل دائمی پرونده هسته‌ای تحت نظارت IAEA.
\end{enumerate}
\end{modernbox}

\newpage
\begin{landscape}
\section{نقشه‌ی راه دیپلماتیک (سال اول گذار)}

\begin{pgfgantt}[
    vgrid, hgrid,
    title/.style={fill=MainBlue, text=white},
    canvas/.style={fill=white},
    group/.style={fill=SlateGray},
    bar/.style={fill=LightBlue},
    milestone/.style={fill=AccentGold},
    y unit chart=0.7cm,
    x unit=0.5cm,
    time slot format=isodate
]{2025-10-01}{2026-10-01}
    \gantttitle{تقویم دیپلماتیک ایران نوین (۲۰۲۵-۲۰۲۶)}{13} \\
    \gantttitlecalendar{year, month=shortname} \\
    
    \ganttbar{پیام به سازمان ملل و IAEA}{2025-10-01}{2025-10-15} \\
    \ganttbar{مذاکرات رفع تحریم}{2025-11-01}{2026-03-31} \\
    \ganttmilestone{لغو قطعنامه‌های تحریمی}{2026-04-01} \\
    \ganttbar{عادی‌سازی روابط منطقه‌ای}{2025-12-01}{2026-09-30} \\
    \ganttbar{الحاق به کنوانسیون‌های حقوق بشر}{2026-01-01}{2026-06-30} \\
    \ganttmilestone{نامزدی شورای حقوق بشر}{2026-10-01}
\end{pgfgantt}
\end{landscape}

\section{مدیریت بحران‌های بین‌المللی}

\begin{table}[htbp]
\centering
\caption{تغییرات کلیدی در سیاست خارجی}
\begin{tabular}{R{4cm} R{5cm} R{6cm}}
\hline
\textbf{موضوع} & \textbf{سیاست قبلی} & \textbf{سیاست پیشنهادی} \\
\hline
برنامه هسته‌ای & تقابل و ابهام & شفافیت تحت \lr{NPT} و \lr{IAEA} \\
روابط منطقه‌ای & صدور انقلاب و نیروهای نیابتی & حسن همجواری و امنیت پیمانی \\
حقوق بشر & مقاومت در برابر پایش & همکاری کامل با سازمان ملل \\
روابط با غرب & «مرگ بر...» و تقابل & تعامل اقتصادی و سیاسی متوازن \\
روابط با شرق & وابستگی راهبردی & تنوع در شراکت بر اساس منافع ملی \\
\hline
\end{tabular}
\end{table}
